\section{Размерности 5 и 7: $\chi(\R^5)\le132$ и $\chi(\R^7)\le1323$}
\label{sec:r57}

В этих размерностях исчерпывающий перебор подрешёток недоступен, и обе
конструкции получены чередованием дискретного поиска ядра с непрерывной
деформацией формы Грама (\S\ref{sec:big}). Ниже сначала описано, как найдено
ядро, затем приведён точный сертификат.

\subsection{Спуск $140\to132$ в $\R^5$}

Прежняя оценка $140$ получена АБПР на решётке $A_5^*$, и там же доказана её
минимальность --- но лишь среди подрешёток \emph{этой} решётки. Наш спуск идёт
ступенями $140\to139\to136\to134\to133\to132$, и на каждой меняются и ядро, и
форма Грама. Первая ступень описана в \S\ref{sec:big} как иллюстрация
примарного шага; опишем остальные.

Для промежуточного результата $136$ проверены типы
$[17,8]$, $[17,4,2]$ и $[17,2,2,2]$. Тип $[17,8]$ после повторного обмена
геометрическими весами и деформации формы пересёк порог в двух независимых
стохастических продолжениях. Доводка по активному набору дала численный минимум
$1{,}0008145796$.

Локальный максимум при $134=67\cdot2$ удалось обойти так. Вместо ширины одного
ядра мы максимизировали вторую по величине из ширин двух ядер сразу: полученная
форма Грама хороша не для одного из них, а для обоих. На этой форме
исчерпывающий перебор с порогом $0{,}90$ дал $28$ новых пригодных ядер --- по
одному на каждую допустимую строку проективного пространства. Четыре лучших из этих ядер были независимо продолжены по метрике;
второе пересекло единицу уже на $97$-м поколении CMA-ES и дало
промежуточную доказанную оценку~$134$.

Затем на полученной форме были намеренно запущены поиски сразу при индексах
$133$, $132$ и $131$. Для $133$ две независимые деформации пересекли единицу.
Для $132$ блочные типы $[11,3,2,2]$ и $[11,3,4]$ дали три пересечения;
лучшее из них доводка по активному набору подняла с $1{,}003354222\ldots$ до
$1{,}010905284\ldots$. Это демонстрирует немонотонность качества соседних
арифметических типов: проверка сразу на два цвета ниже оказалась успешнее
остановки на первом пересечённом индексе. Именно конструкция индекса~$132$
рационализована ниже.

\begin{theorem}\label{thm:r5}
Существует решёточная раскраска
\[
 \chi\bigl(\R^5,[1,\ell]\bigr)\le132
 \qquad\text{для всех }\ell\le\frac{101}{100}.
\]
В частности, $\chi(\R^5)\le132$.
\end{theorem}

\begin{proof}[Компьютерно проверяемое доказательство]
Пусть $\Lambda=B^{\mathsf T}\Z^5$, где $BB^{\mathsf T}=Q$ и
\[
Q=\frac1{100000}
\begin{pmatrix}
79327&86284&69597&36140&26501\\
86284&254538&216808&101257&102135\\
69597&216808&296568&163667&140527\\
36140&101257&163667&181733&71653\\
26501&102135&140527&71653&148076
\end{pmatrix}.
\]
Её главные угловые миноры
\[
79327,\ 12746807270,\ 1422467480626358,\
124437401434750889345,\ 9999935596575703407615413
\]
положительны. Зададим цвет координатного вектора $x\in\Z^5$ тройкой
гомоморфизмов
\[
\begin{aligned}
\varphi_{11}(x)&=x_2+6x_3+2x_4+3x_5\pmod {11},\\
\varphi_3(x)&=x_1+2x_2+x_3\pmod 3,\\
\varphi_4(x)&=x_2+3x_3+x_4+3x_5\pmod 4.
\end{aligned}
\]
По китайской теореме об остатках это эквивалентно одной строке
\[
\varphi(x)=88x_1+89x_2+127x_3+57x_4+3x_5\pmod {132}.
\]
Положим $\Gamma=B^{\mathsf T}\ker\varphi$. Гомоморфизм сюръективен,
$[\Lambda:\Gamma]=132$, а Smith-диагональ ядра равна
$(1,1,1,1,132)$.

Ячейка Вороного рациональной формы $Q$ имеет $62$ фасеты и $720$ вершин.
Все вершинные системы решены над $\mathbb Q$; точный квадрат радиуса покрытия
равен
\[
R^2=
\frac{3093570095915736710467707313641}
     {3999974238630281363046165200000}.
\]
Доказываемое утверждение --- это $D(v)\ge\ell\,\diam V_0$ при $\ell=101/100$,
поэтому из $D(v)\ge\|v\|-\diam V_0=\|v\|-2R$ следует, что возможный нарушитель
обязан удовлетворять
\begin{equation}\label{eq:win5}
\|v\|<2(1+\ell)R,\qquad\text{т.е.}\qquad \|v\|^2<12{,}4985\ldots
\end{equation}
Отметим, что более узкое окно $\|v\|<4R$ отвечает лишь цели $d>1$ и для
интервального утверждения недостаточно. Перебор Финке--Поста по LLL-приведённому
базису $\Gamma$ оставляет в окне~\eqref{eq:win5} ровно $38$ ненулевых векторов
($19$ пар); из них $36$ лежат уже в окне $\|v\|<4R$ и независимо воспроизводятся
C\raise.1ex\hbox{\small++}-перечислителем, а оставшаяся пара
$\pm(-2,2,-2,2,2)$ даёт $D/\diam V_0=1{,}0442\ldots$ и на минимум не влияет.

Для всех $38$ векторов задача проекции на $V_0$ решена точно: допустимость,
неотрицательность множителей ККТ и значения целей проверены над $\mathbb Q$.
Минимум достигается на паре
$\pm(2,0,1,-3,0)$ (координаты по базису $\Lambda$) и равен
\[
\begin{aligned}
D_{\min}^2&=
\frac{1634114770527727}{516898614620000},\\
D_{\min}^2-\left(\frac{101}{100}\right)^2 4R^2
&=\frac{1450500657237822255902962097875907780124229}
{258447642807960215367421006872956903000000000}>0.
\end{aligned}
\]
Следовательно,
$D_{\min}/\diam V_0=1{,}010897714548927\ldots>1{,}01$.
Раскраска ячеек по классам $\Lambda/\Gamma$ даёт требуемые $132$ цвета.
Полный список вершинных инцидентностей и KKT-сертификатов, а также независимая
проверка множества из $1025$ пар $\pm v$, запрещённых при пороге $\ell=1$,
находятся в файле
\path{audit-data/results/metric_deform_a5_132_refined_certificate.json}.
Дополнительный верификатор, не использующий Qhull или вещественный
перечислитель коротких векторов, точно проверяет каждый из $31$ ненулевого
класса чётности, строит $62$ фасеты, обходит связный граф из $720$ вершин и
$1800$ рёбер и воспроизводит те же $R^2$, $D_{\min}^2$ и KKT-сертификаты;
его протокол ---
\path{audit-data/results/metric_deform_a5_132_refined_independent_exact_audit.json}.
Наконец, весь независимый аудит повторён после рационализации той же численной
формы со знаменателями $10^4$ и $10^6$; получены отношения соответственно
$1{,}010770760\ldots$ и $1{,}010904642\ldots$, и в обоих случаях также
сохраняется строгий запас для $\ell=1{,}01$.
\end{proof}

\paragraph{Барьер $1372$ на фиксированной решётке $E_7^*$.} Для $\R^7$
\cite{ABPR} получили оценку $1372$ решёткой $E_7^*$ \emph{рандомизированным}
поиском, без доказательства минимальности (в отличие от $\R^5$, где минимальность
$140$ доказана полным перебором). Мы воспроизвели их конструкцию точно (матрицы
$M$, $C_7$ из \cite{ABPR}: $|F_1|=24\,284$, $|\det C_7|=1372$, $\Lambda'\cap
F=\varnothing$; отношение покрытие/упаковка $E_7^*$ равно $\sqrt{7/3}$) и
проверили её на понижение индекса. Подрешётка $C_7$ оказалась
\emph{неуплотняемой}: не существует \emph{надрешётки} $\Lambda''\supset\Lambda'$
индекса $1372/p$ ($p\in\{2,7\}$), избегающей $F$ (любой склеивающий вектор
порядка $p$ попадает в $F$). Ни градиентный спуск, ни отжиг, ни $2\cdot10^6$
случайных базисов из коротких незапрещённых векторов не дали пригодной подрешётки
индекса ${<}1372$ (случайные пригодные подрешётки встречаются лишь с индексом
${\ge}2052$). Это --- вычислительное свидетельство оптимальности $1372$ среди
раскрасок с фиксированной родительской решёткой $E_7^*$, но не среди всех форм
Грама.


\subsection{Спуск $1372\to1323$ в $\R^7$}

Прямые атаки на фиксированную $E_7^*$ упирались в барьер (см.\ выше), поэтому
здесь понадобились обе стадии сразу: примарный Radon-поиск ядра
(\S\ref{sec:big}) и последующая деформация формы Грама. При структуре
$(\Z/7)^2\times(\Z/3)^3$ поиск дал ядро индекса $1323$, у которого на
$E_7^*$ остаются лишь шесть конфликтов одинаковой тяжести
$D/\diam V_0=0{,}886405\ldots$. После этого ядро фиксируется, а родительский
базис деформируется как $B=B_0\exp H$, где $H=H^{\mathsf T}$ и
$\operatorname{tr}H=0$. В отличие от прежней целочисленной цели --- числа неотделённых векторов ---
непрерывной целью служит непосредственно $\min D/\diam V_0$.

\begin{theorem}\label{thm:r7}
Существует решёточная раскраска
\[
 \chi\bigl(\R^7,[1,\ell]\bigr)\le1323
 \qquad\text{для всех }\ell\le1{,}007.
\]
В частности, $\chi(\R^7)\le1323$.
\end{theorem}

\begin{proof}[Компьютерно проверяемое доказательство]
Пусть $\Lambda=B^{\mathsf T}\Z^7$, где $BB^{\mathsf T}=Q$ и
\[
\begingroup
\setlength{\arraycolsep}{2.6pt}
Q=\frac1{10\,000}\!
\begin{pmatrix}
20968&-10761&-529&-2493&6504&-6603&-796\\
-10761&22451&-9533&1398&-2933&1926&13148\\
-529&-9533&21640&-12309&-430&1476&-630\\
-2493&1398&-12309&22527&-9467&974&1288\\
6504&-2933&-430&-9467&18435&-11881&-1083\\
-6603&1926&1476&974&-11881&20755&-99\\
-796&13148&-630&1288&-1083&-99&18191
\end{pmatrix}.
\endgroup
\]
Все главные угловые миноры числителя положительны. Зададим
$\Gamma=B^{\mathsf T}C\Z^7$ матрицей
\[
\begingroup
\setlength{\arraycolsep}{3.2pt}
C=\begin{pmatrix}
21&0&9&17&16&14&2\\
0&21&9&0&10&14&20\\
0&0&3&1&1&2&0\\
0&0&0&1&0&0&0\\
0&0&0&0&1&0&0\\
0&0&0&0&0&1&0\\
0&0&0&0&0&0&1
\end{pmatrix}.
\endgroup
\]
Имеем $|\det C|=1323$, а нормальная форма Смита равна
$(1,1,1,1,3,21,21)$.

Для рациональной формы $Q$ ячейка Вороного имеет $254$ фасеты и является
простым многогранником с $39\,296$ вершинами. Все вершинные системы решены над
$\mathbb Q$; точный квадрат радиуса покрытия равен
\[
 R^2=
 \frac{42823481847474795123735510069131}
 {49956306567646383875936777960000}.
\]
Доказываемое утверждение --- это $D(v)\ge\ell\,\diam V_0$ при $\ell=1{,}007$,
поэтому из $D(v)\ge\|v\|-\diam V_0=\|v\|-2R$ следует, что возможный нарушитель
обязан удовлетворять
\begin{equation}\label{eq:win7}
 \|v\|<2(1+\ell)R,\qquad\text{т.е.}\qquad\|v\|^2<13{,}8117\ldots
\end{equation}
(более узкое окно $\|v\|<4R$ отвечало бы лишь цели $d>1$). Перебор Финке--Поста
по LLL-приведённому базису $\Gamma$ оставляет в окне~\eqref{eq:win7} ровно $72$
ненулевых вектора ($36$ пар); из них $70$ лежат уже в окне $\|v\|<4R$, а
оставшаяся пара $\pm(2,-1,0,0,0,0,1)$ даёт $D/\diam V_0=1{,}18905\ldots$ и на
минимум не влияет. Для каждого из них решена задача проекции на $V_0$,
причём допустимость, неотрицательность множителей ККТ и значение цели проверены
над $\mathbb Q$. Минимум достигается на паре
$\pm(1,0,0,2,1,0,1)$ (координаты по базису $\Lambda$) и равен
\[
 D_{\min}^2=\frac{21004929056}{6040637375},\qquad
 D_{\min}^2-\left(\tfrac{1007}{1000}\right)^2 4R^2>0,
\]
где последняя разность равна в точности
\[
 \frac{134661904230636621379975888187452316904519}
 {603535865138965424449162128164504510000000000}.
\]
Следовательно,
$D_{\min}/\diam V_0=1{,}007032309005\ldots>1{,}007$.
Раскрашивая ячейки Вороного по классам $\Lambda/\Gamma$, получаем требуемые
$1323$ цвета. Счётчики вершин и фасет, сертификаты ККТ и независимая проверка
$\ker\varphi\cap F=\varnothing$ находятся в файле
\path{audit-data/results/metric_deform_e7_1323_certificate.json}.
\end{proof}

\paragraph{Границы \emph{клеточных} раскрасок на фиксированной форме.}
Схема АБПР --- частный случай \emph{клеточной} раскраски (цвет
постоянен на внутренности ячейки Вороного). Общая клеточная раскраска с периодом
$\Gamma$ --- это правильная раскраска \emph{графа конфликтов} $H_\Gamma$ на
$\Lambda/\Gamma$ (метрический аналог раскраски графа касания ячеек Вороного
\cite{DSMMV}): вершины --- классы, ребро $[a]\!\sim\![b]$, если пара ячеек
реализует запрещённое расстояние. АБПР дают каждому классу свой цвет
($[\Lambda:\Lambda']$ цветов), но $\chi(H_\Gamma)$ \emph{в принципе} может быть
меньше. Мы установили, что здесь это не так.

\emph{(i) Граф не разрежается.} Множество рёбер \emph{точно} отвечает $F$:
для всякого $v\in F$ пара ячеек $V_0$ и $v+V_0$ действительно реализует
запрещённое расстояние, поскольку наряду с $D(v)<\diam V_0$ выполнено и
$D'(v)\ge2R$ для наибольшего расстояния $D'(v)$ между этими ячейками (проверено
перебором вершин): функция расстояния непрерывна на связном множестве
$V_0\times(v+V_0)$, поэтому реализуются \emph{все} значения из
$[D(v),D'(v)]$, в частности запрещённое $2R$.

\emph{(ii) При периоде $\Gamma=\Lambda'$ выигрыша нет.} Здесь $F$ покрывает
\emph{все} ненулевые классы $\Lambda/\Lambda'$ (проверено для $A_5^*/140$ и
$E_7^*/1372$), так что $H_{\Lambda'}$ --- полный граф
$K_{[\Lambda:\Lambda']}$ и $\chi(H_{\Lambda'})=[\Lambda:\Lambda']$.

\emph{(iii) Измельчение периода тоже не помогает.} Для проверенных
$\Gamma\subsetneq\Lambda'$ граф конфликтов оказывается полным многодольным; в
$\R^5$ при периоде индекса $280$ равенство $\chi(H_\Gamma)=140$ \emph{доказано}
SAT-решателем. Клика у этих измельчённых графов меньше индекса ($80$ в $\R^5$
и $541$ в $\R^7$), так что чисто кликовая нижняя оценка здесь недостаточна, --- и
тем не менее во всех опробованных случаях хроматическое число совпадало с
индексом. Эти отрицательные результаты относятся к фиксированным
симметричным родительским решёткам. Теорема~\ref{thm:r7} показывает, что
совместная смена ядра и формы Грама обходит барьер без перехода к
подъячеечным раскраскам. Строгое утверждение о клеточных раскрасках на
фиксированной $E_7^*$ остаётся открытым.


